% GENERATED FILE. Edit canonical lessons, then run npm run generate:books.
% canonical-source-hash: fnv1a64:eedb4dc1248f3447
% canonical-lessons: PT-C522-exposicao, PT-C522-galeria, PT-C522-retrato, PT-C522-esculpir, PT-C522-escultor, PT-R522-revisao-515-522

\chapter{Na galeria}
\label{ch:pt-na-galeria}
\hlchaptermodality{522}

\begin{chapteropening}
\textbf{By the end of this chapter:} \emph{I can talk about exhibitions, portraits and sculpture.}

Complete the last lesson of chapter 522: I can talk about exhibitions, portraits and sculpture.
\end{chapteropening}

\section[a exposição]{a exposição — exhibition}

\label{lesson:PT-C522-exposicao}



\begin{quote}
Before the new one: say the Portuguese for ballerina, then the Portuguese for (theatre) dressing room.
\end{quote}

\subsection*{You'll want to know: a exposição}
\textbf{a exposição} — \textquotedblleft{}exhibition\textquotedblright{}.

\begin{quote}
Há uma exposição de fotografia no museu.
\end{quote}

\begin{quote}
\emph{There is a photography exhibition at the museum.}
\end{quote}

\emph{Inaugurar uma exposição} — to open an exhibition.

\begin{cousinweb}
\emph{exposição} is a learned borrowing from Latin \emph{expositiō} \textquotedblleft{}a setting out\textquotedblright{} (EtymDB), from \emph{expōnere}, built on \emph{pōnere} \textquotedblleft{}to put\textquotedblright{}.
\end{cousinweb}

\begin{grammarlens}[title={What you've built}]
One more word for telling, explaining and describing at B1.
\end{grammarlens}

\subsection*{Your turn}
\begin{itemize}
\raggedright
  \item \emph{Say it:} \emph{a exposição}
  \item \emph{Say it:} \emph{a exposição}, once more
  \item \emph{Say it:} a sentence of your own with \emph{a exposição}
  \item \emph{Recall:} say the Portuguese for ballerina, then the Portuguese for (theatre) dressing room, then say \emph{a exposição} again
\end{itemize}

\begin{tcolorbox}[breakable,colback=teal!4,colframe=teal!35!black,title={Before you move on}]
What does a exposição mean? (\textquotedblleft{}Exhibition\textquotedblright{}.) Say it once more.
\end{tcolorbox}
\section[a galeria]{a galeria — gallery}

\label{lesson:PT-C522-galeria}



\begin{quote}
Before the new one: say the Portuguese for (theatre) dressing room, then the Portuguese for exhibition.
\end{quote}

\subsection*{You'll want to know: a galeria}
\textbf{a galeria} — \textquotedblleft{}gallery\textquotedblright{}.

\begin{quote}
A galeria fica perto do rio.
\end{quote}

\begin{quote}
\emph{The gallery is near the river.}
\end{quote}

\emph{Galeria de arte} — art gallery.

\begin{cousinweb}
\emph{galeria} was borrowed from Italian \emph{galleria} (EtymDB), from Medieval Latin \emph{galeria}, whose deeper origin is not settled.
\end{cousinweb}

\begin{grammarlens}[title={What you've built}]
One more word for telling, explaining and describing at B1.
\end{grammarlens}

\subsection*{Your turn}
\begin{itemize}
\raggedright
  \item \emph{Say it:} \emph{a galeria}
  \item \emph{Say it:} \emph{a galeria}, once more
  \item \emph{Say it:} a sentence of your own with \emph{a galeria}
  \item \emph{Recall:} say the Portuguese for (theatre) dressing room, then the Portuguese for exhibition, then say \emph{a galeria} again
\end{itemize}

\begin{tcolorbox}[breakable,colback=teal!4,colframe=teal!35!black,title={Before you move on}]
What does a galeria mean? (\textquotedblleft{}Gallery\textquotedblright{}.) Say it once more.
\end{tcolorbox}
\section[o retrato]{o retrato — portrait}

\label{lesson:PT-C522-retrato}



\begin{quote}
Before the new one: say the Portuguese for exhibition, then the Portuguese for gallery.
\end{quote}

\subsection*{You'll want to know: o retrato}
\textbf{o retrato} — \textquotedblleft{}portrait\textquotedblright{}.

\begin{quote}
Na sala há um retrato da bisavó.
\end{quote}

\begin{quote}
\emph{In the living room there is a portrait of my great-grandmother.}
\end{quote}

Also a photo of a person.

\begin{cousinweb}
\emph{retrato} comes from Italian \emph{ritratto} (EtymDB), from Latin \emph{retractus}, built on \emph{retrahere} \textquotedblleft{}to draw back\textquotedblright{}, from \emph{trahere} \textquotedblleft{}to draw\textquotedblright{}.
\end{cousinweb}

\begin{grammarlens}[title={What you've built}]
One more word for telling, explaining and describing at B1.
\end{grammarlens}

\subsection*{Your turn}
\begin{itemize}
\raggedright
  \item \emph{Say it:} \emph{o retrato}
  \item \emph{Say it:} \emph{o retrato}, once more
  \item \emph{Say it:} a sentence of your own with \emph{o retrato}
  \item \emph{Recall:} say the Portuguese for exhibition, then the Portuguese for gallery, then say \emph{o retrato} again
\end{itemize}

\begin{tcolorbox}[breakable,colback=teal!4,colframe=teal!35!black,title={Before you move on}]
What does o retrato mean? (\textquotedblleft{}Portrait\textquotedblright{}.) Say it once more.
\end{tcolorbox}
\section[esculpir]{esculpir — to sculpt, to carve}

\label{lesson:PT-C522-esculpir}



\begin{quote}
Before the new one: say the Portuguese for gallery, then the Portuguese for portrait.
\end{quote}

\subsection*{You'll want to know: esculpir}
\textbf{esculpir} — \textquotedblleft{}to sculpt, to carve\textquotedblright{}.

\begin{quote}
O artista esculpiu a figura em pedra.
\end{quote}

\begin{quote}
\emph{The artist carved the figure in stone.}
\end{quote}

Sculpture is \emph{a escultura}.

\begin{cousinweb}
\emph{esculpir} is a learned borrowing from Latin \emph{sculpere} \textquotedblleft{}to carve, to chisel\textquotedblright{} (Lewis \& Short).
\end{cousinweb}

\begin{grammarlens}[title={What you've built}]
One more word for telling, explaining and describing at B1.
\end{grammarlens}

\subsection*{Your turn}
\begin{itemize}
\raggedright
  \item \emph{Say it:} \emph{esculpir}
  \item \emph{Say it:} \emph{esculpir}, once more
  \item \emph{Say it:} a sentence of your own with \emph{esculpir}
  \item \emph{Recall:} say the Portuguese for gallery, then the Portuguese for portrait, then say \emph{esculpir} again
\end{itemize}

\begin{tcolorbox}[breakable,colback=teal!4,colframe=teal!35!black,title={Before you move on}]
What does esculpir mean? (\textquotedblleft{}To sculpt, to carve\textquotedblright{}.) Say it once more.
\end{tcolorbox}
\section[o escultor]{o escultor — sculptor}

\label{lesson:PT-C522-escultor}



\begin{quote}
Before the new one: say the Portuguese for portrait, then the Portuguese for to sculpt.
\end{quote}

\subsection*{You'll want to know: o escultor}
\textbf{o escultor} — \textquotedblleft{}sculptor\textquotedblright{}.

\begin{quote}
O escultor trabalha com madeira.
\end{quote}

\begin{quote}
\emph{The sculptor works with wood.}
\end{quote}

A woman is \emph{a escultora}.

\begin{cousinweb}
\emph{escultor} comes from Latin \emph{sculptor} \textquotedblleft{}a carver\textquotedblright{} (EtymDB), from \emph{sculpere} \textquotedblleft{}to carve\textquotedblright{}.
\end{cousinweb}

\begin{grammarlens}[title={What you've built}]
One more word for telling, explaining and describing at B1.
\end{grammarlens}

\subsection*{Your turn}
\begin{itemize}
\raggedright
  \item \emph{Say it:} \emph{o escultor}
  \item \emph{Say it:} \emph{o escultor}, once more
  \item \emph{Say it:} a sentence of your own with \emph{o escultor}
  \item \emph{Recall:} say the Portuguese for portrait, then the Portuguese for to sculpt, then say \emph{o escultor} again
\end{itemize}

\begin{tcolorbox}[breakable,colback=teal!4,colframe=teal!35!black,title={Before you move on}]
What does o escultor mean? (\textquotedblleft{}Sculptor\textquotedblright{}.) Say it once more.
\end{tcolorbox}
\section[(dialogue)]{Revisão dos capítulos 515–522}

\label{lesson:PT-R522-revisao-515-522}



\begin{quote}
Say the Portuguese for changing room and supporters' group.
\end{quote}

\subsection*{Your turn}
Say these aloud:

\begin{itemize}
\raggedright
  \item changing room, supporters' group, injury, to get injured, cheating
  \item orchestra, conductor, symphony, cello, harp
  \item flute, trumpet, saxophone, accordion, to tune
  \item choir, melody, chorus, to hum, to whistle
  \item rehearsal, to rehearse, stage, audience, applause
  \item set, to stage, staging, premiere, to premiere
  \item comedy, opera, ballet, ballerina, (theatre) dressing room
  \item exhibition, gallery, portrait, to sculpt, sculptor
\end{itemize}

\begin{tcolorbox}[breakable,colback=teal!4,colframe=teal!35!black,title={Before you move on}]
Changing room? (\textbf{o balneário}.) Supporters' group? (\textbf{a claque}.) Injury? (\textbf{a lesão}.) To get injured? (\textbf{lesionar}.) Cheating? (\textbf{a batota}.) Orchestra? (\textbf{a orquestra}.) Conductor? (\textbf{o maestro}.) Symphony? (\textbf{a sinfonia}.) Cello? (\textbf{o violoncelo}.) Harp? (\textbf{a harpa}.) Flute? (\textbf{a flauta}.) Trumpet? (\textbf{o trompete}.) Saxophone? (\textbf{o saxofone}.) Accordion? (\textbf{o acordeão}.) To tune? (\textbf{afinar}.) Choir? (\textbf{o coro}.) Melody? (\textbf{a melodia}.) Chorus? (\textbf{o refrão}.) To hum? (\textbf{cantarolar}.) To whistle? (\textbf{assobiar}.) Rehearsal? (\textbf{o ensaio}.) To rehearse? (\textbf{ensaiar}.) Stage? (\textbf{o palco}.) Audience? (\textbf{a plateia}.) Applause? (\textbf{o aplauso}.) Set? (\textbf{o cenário}.) To stage? (\textbf{encenar}.) Staging? (\textbf{a encenação}.) Premiere? (\textbf{a estreia}.) To premiere? (\textbf{estrear}.) Comedy? (\textbf{a comédia}.) Opera? (\textbf{a ópera}.) Ballet? (\textbf{o bailado}.) Ballerina? (\textbf{a bailarina}.) (theatre) dressing room? (\textbf{o camarim}.) Exhibition? (\textbf{a exposição}.) Gallery? (\textbf{a galeria}.) Portrait? (\textbf{o retrato}.) To sculpt? (\textbf{esculpir}.) Sculptor? (\textbf{o escultor}.)
\end{tcolorbox}
